%============================================================ 6. 결과
\section{결과}
\label{sec:results}

\ref{sec:experiment}절의 $2\times2$ 설계는 두 계층의 기여를 따로 잰다. 이 절은 그 결과를
기동 계층 단독의 효과(\ref{sec:res_g1}절), 지휘관 계층의 효과와 그것이 성립하는
조건(\ref{sec:res_g2}절), 두 효과의 결합 방식(\ref{sec:res_g3}절) 순으로 보고하고, 배치
관점의 지연 측정을 덧붙인다(\ref{sec:res_latency}절).

\subsection{기동 계층의 효과: 자원과 기동의 절감}
\label{sec:res_g1}

\subsubsection{기동 계층: 휴리스틱 대 U-Net 점수맵}
\label{sec:res_maneuver}

언어모델 없이 \ref{sec:heuristic}절의 기하 휴리스틱과 \ref{sec:policy}절의 점수맵 정책만
비교한다. 배정은 둘 다 휴리스틱이므로 차이는 기동 방식뿐이다(Table~\ref{tab:maneuver}). 3개 포메이션 통합 기준으로 포획률은 $0.844 \to 0.886$
($+0.041$, $\dz = 0.37$)로 유의하게 개선되었다.
다만 지표 10개에 대한 다중비교 보정을 걸면 구간이 $[-0.013, +0.100]$ 로 0을
포함하여(\ref{sec:stats}절) 이 개선은 유의성을 잃는다. 이 절의 요점은 포획률의 증가가
아니라 \emph{같은 정도로 막으면서 자원과 기동을 크게 줄인다}는 데 있다.

% 자동 생성 — tools/make_paper_tables.py. 손으로 고치지 말 것.
\begin{table*}[tp]
\centering
\caption{기동 계층(U-Net 점수맵 + GRPO)의 단독 기여. 두 조건 모두 배정은 휴리스틱이며 \textbf{LLM 을 전혀 사용하지 않는다}. 기여는 파상 포메이션에서 가장 크다.}
\label{tab:maneuver}
\footnotesize
\resizebox{\ifdim\width>\linewidth\linewidth\else\width\fi}{!}{%
\begin{tabular}{lcccc}
\toprule
포메이션 & 휴리스틱 기동 & U-Net 기동 & $\Delta$ [95\,\% CI] & $d_z$ \\
\midrule
집중 & 0.997 & 0.997 & +0.000 [-0.010, +0.010] & +0.00 \\
양동 & 0.787 & 0.753 & -0.033 [-0.077, +0.000] & -0.31 \\
파상 & 0.750 & 0.907 & +0.157 [+0.043, +0.270]$^{*}$ & +0.48 \\
\textbf{전체} & 0.844 & 0.886 & +0.041 [+0.003, +0.082]$^{*}$ & +0.37 \\
\bottomrule
\end{tabular}
}
\\[3pt]
\begin{minipage}{\linewidth}\footnotesize
$^{*}$ 신뢰구간이 0 을 포함하지 않음. 전체 행은 3개 포메이션 통합.
\end{minipage}
\end{table*}


\begin{figure}[!tbp]
\centering
\includegraphics[width=\columnwidth]{fig15.png}
\caption{휴리스틱 대 U-Net 기동의 지표 프로파일(바깥일수록 우수): (a) 전체, (b) 집중,
         (c) 양동, (d) 파상}
\label{fig:radar_maneuver}
\end{figure}

같은 비교를 10개 지표 전부에 걸쳐 보인 것이 Fig.~\ref{fig:radar_maneuver}다. 각 축은 8개
조건의 범위로 0--1 정규화하여 바깥일수록 우수하다. U-Net 은 포획당 이동거리 축에서만
안쪽이다 --- 자리를 골라 그물벽을 놓는 대신 더 멀리 움직인다. 충돌률·전개 시 적 거리 축은
겹치고, 나머지 일곱 축에서 바깥에 있다.

포메이션별로 보면 기여가 고르지 않다. 파상 포메이션에서 $+0.157$($\dz = 0.48$)로 가장 크고,
집중 포메이션에서는 $0.000$, 양동 포메이션에서는 $-0.033$ 으로 오히려 소폭 낮다. 다만 후자의
차이가 유의하지 않아 악화라고 말할 수 없다.
파상의 $+0.157$ 이 소수 시드가 끄는 값은 아니다. 시드 대응 차이의 분포(Fig.~\ref{fig:paired_diff})를
보면 30개 시드 가운데 15개에서 양, 12개에서 0, 3개에서 음이고 중앙값은 $+0.050$ 이다.
양동에서는 양 2개, 음 10개로 방향이 반대다. 이 대비가 다음 설명의 근거다.

\begin{figure*}[!tbp]
\centering
\includegraphics[width=0.8\textwidth]{fig16.png}
\caption{시드 대응 차이의 분포(포획률, 기준선 heur\_heur 대비). 점은 시드, 마름모는 평균과 BCa 95\,\%
         신뢰구간; 왼쪽 열은 기동 계층 단독(heur\_unet), 오른쪽 열은 제안 체계(llm\_unet, Gemini).}
\label{fig:paired_diff}
\end{figure*}

이 분포는 \ref{sec:scenario}절의 포메이션 구조와 대응한다. 파상에서는 앞 단을 막은 그물벽의
\emph{위치}가 뒤 단의 차단을 좌우하므로 해면 전체를 평가하는 점수맵이 직접 효과를 내고, 양동에서는
무리 수와 방어정 수가 같아 결정적인 것이 \emph{배정}이라 배정이 어긋나면 기동이 보정하지
못한다. 집중은 \ref{sec:res_ceiling}절에서 보이듯 천장 구간이다.

\subsubsection{학습 알고리즘: 같은 행동공간의 그룹 상대 최적화와 MAPPO}
\label{sec:res_algo}

점수맵 정책의 이득이 픽셀 지목이라는 표현에서 오는지 가치망 없는 갱신 규칙에서 오는지를
가르기 위해, 같은 액터·행동공간·보상·예산 위에서 갱신 규칙만 MAPPO 로 바꾼 정책 세 개를 같은
프로토콜로 평가하였다(Table~\ref{tab:algo}, Fig.~\ref{fig:algo}).

% 자동 생성 — tools/make_paper_tables.py. 손으로 고치지 말 것.
\begin{table*}[t]
\centering
\caption{같은 픽셀 행동공간 위의 학습 알고리즘 비교. 배정은 모두 휴리스틱이며 기동 정책의 액터·보상·초기화·환경 step 예산이 같고 갱신 규칙만 다르다. MAPPO 는 학습 시드 3런의 에피소드별 평균, $\Delta$ 는 GRPO $-$ MAPPO 의 시드 대응 차이.}
\label{tab:algo}
\footnotesize
\resizebox{\ifdim\width>\linewidth\linewidth\else\width\fi}{!}{%
\begin{tabular}{lccccc}
\toprule
지표 & 휴리스틱 & MAPPO & GRPO(채택) & $\Delta$ [95\,\% CI] & $d_z$ \\
\midrule
포획률 $\uparrow$ & 0.844 & 0.859 & 0.886 & +0.026 [-0.003, +0.056] & +0.31 \\
돌파 수 $\downarrow$ & 1.56 & 1.41 & 1.14 & -0.26 [-0.57, +0.03] & -0.31 \\
충돌률(척당) $\downarrow$ & 0.130 & 0.138 & 0.119 & -0.020 [-0.047, +0.004] & -0.27 \\
그물 걸림 $\downarrow$ & 0.82 & 0.70 & 0.62 & -0.08 [-0.23, +0.05] & -0.20 \\
포획당 그물 $\downarrow$ & 0.737 & 0.733 & 0.568 & -0.165 [-0.213, -0.116]$^{*}$$^{\dagger}$ & -1.19 \\
포획당 이동거리 (m) $\downarrow$ & 825 & 813 & 896 & +82 [+40, +139]$^{*}$$^{\dagger}$ & +0.59 \\
총 선회량 (rad) $\downarrow$ & 1625 & 1541 & 1464 & -77 [-158, +17] & -0.31 \\
포획 이격거리 (m) $\uparrow$ & 1511 & 1523 & 1678 & +154 [+112, +202]$^{*}$$^{\dagger}$ & +1.21 \\
평균 포획 시각 (step) $\downarrow$ & 430 & 430 & 418 & -12 [-19, -5]$^{*}$$^{\dagger}$ & -0.58 \\
전개 시 적 거리 (m) $\downarrow$ & 1499 & 1525 & 1425 & -100 [-154, -51]$^{*}$$^{\dagger}$ & -0.68 \\
\bottomrule
\end{tabular}
}
\\[3pt]
\begin{minipage}{\linewidth}\footnotesize
$^{*}$ 95\,\% 구간이 0 을 배제. $^{\dagger}$ Bonferroni 보정(99.5\,\%)에서도 배제. MAPPO 런별 포획률 0.876, 0.838, 0.864. 3개 포메이션 통합(각 조건 30시드).
\end{minipage}
\end{table*}


\begin{figure}[!tbp]
\centering
\includegraphics[width=0.92\columnwidth]{fig17.png}
\caption{같은 픽셀 행동공간 위의 휴리스틱·MAPPO·GRPO 기동의 지표 프로파일(바깥일수록 우수)}
\label{fig:algo}
\end{figure}

포획률은 갈리지 않는다. MAPPO $0.859$ 와 GRPO $0.886$ 의 차이는 유의하지 않고, 포메이션별
방향도 파상 개선·양동 저하로 같다. 갈리는 것은 \emph{자원과 조기 제압}이다. MAPPO 의
프로파일은 포획률·돌파 축을 빼면 휴리스틱 위에 그대로 겹치는 반면, GRPO 는 포획당 그물을
$0.568$ 로 줄이고($\dz = -1.19$), 포획 이격거리를 $154\unit{m}$ 늘리며($\dz = +1.21$), 평균
포획 시각과 전개 시 적 거리를 각각 12 step 과 $100\unit{m}$ 앞당긴다. 네 지표 모두
다중비교 보정에서도 유의하다. MAPPO 가 이기는 축은 포획당 이동거리 하나다.

즉 같은 표현을 공유하는 한 갱신 규칙은 포획률을 가르지 못하고, 그물 절약과 조기 제압만
갱신 규칙에 좌우된다. 휴리스틱 후보를 기준으로 삼는 그룹 상대 이득(\ref{sec:training}절)은
휴리스틱을 넘어서는 선택에만 양의 이득을 주어 휴리스틱의 자원 사용 패턴에서 벗어나는 반면,
팀 가치로 학습하는 MAPPO 는 어느 배가 그물을 아꼈는지를 개별 배에 귀속시킬 신호가 약해 그
패턴을 물려받는 것으로 해석된다.

\subsection{지휘관 계층의 효과와 성립 조건}
\label{sec:res_g2}

\subsubsection{전체 조건 비교}
\label{sec:res_main}

Table~\ref{tab:main}과 Fig.~\ref{fig:condbars}에 8개 조건의 포획률을 정리하였다. 그림의
무채색 두 막대가 언어모델 없는 기준선이고 그 오른쪽은 모델 규모 오름차순이다.

% 자동 생성 — tools/make_paper_tables.py. 손으로 고치지 말 것.
\begin{table*}[tp]
\centering
\caption{조건별 포획률. 평균과 부트스트랩 95\,\% 신뢰구간, baseline 대비 대응표본 차이. 3개 포메이션 통합(각 조건 30시드).}
\label{tab:main}
\footnotesize
\setlength{\tabcolsep}{4pt}%
\resizebox{\ifdim\width>\linewidth\linewidth\else\width\fi}{!}{%
\begin{tabular}{llccc}
\toprule
지휘관(배정) & 기동 & 포획률 [95\,\% CI] & $\Delta$ vs baseline & $d_z$ \\
\midrule
휴리스틱 (baseline) & 휴리스틱 & 0.844 [0.808, 0.880] & -- & -- \\
휴리스틱 & U-Net 점수맵 & 0.886 [0.847, 0.918] & +0.041 [+0.003, +0.082]$^{*}$ & +0.37 \\
Qwen2.5 7B & 휴리스틱 & 0.652 [0.571, 0.716] & -0.192 [-0.264, -0.129]$^{*}$ & -1.00 \\
Qwen2.5 14B & 휴리스틱 & 0.811 [0.766, 0.854] & -0.033 [-0.086, +0.017] & -0.23 \\
Gemini 3.5 Flash-Lite & 휴리스틱 & 0.904 [0.864, 0.937] & +0.060 [+0.024, +0.100]$^{*}$ & +0.56 \\
Qwen2.5 7B & U-Net 점수맵 & 0.641 [0.582, 0.693] & -0.203 [-0.274, -0.138]$^{*}$ & -1.03 \\
Qwen2.5 14B & U-Net 점수맵 & 0.754 [0.698, 0.806] & -0.090 [-0.153, -0.033]$^{*}$ & -0.54 \\
Gemini 3.5 Flash-Lite & U-Net 점수맵 & 0.907 [0.869, 0.937] & +0.062 [+0.013, +0.113]$^{*}$ & +0.44 \\
\bottomrule
\end{tabular}
}
\\[3pt]
\begin{minipage}{\linewidth}\footnotesize
$^{*}$ 신뢰구간이 0 을 포함하지 않음. $d_z$ 는 대응표본 효과크기. 첫 두 열이 $2\times2$ 설계의 두 축이다. 같은 지휘관 행 두 개를 비교하면 기동 계층의 기여가, 같은 기동 열을 비교하면 지휘관의 기여가 나온다. Qwen2.5 는 로컬(단일 RTX 4070) 서빙, Gemini 는 클라우드 API 다.
\end{minipage}
\end{table*}


\begin{figure*}[!tbp]
\centering
\includegraphics[width=0.8\textwidth]{fig18.png}
\caption{포메이션별·조건별 포획률(색: 지휘관, 빗금: 기동 계층, 오차막대: 시드 부트스트랩 신뢰구간)}
\label{fig:condbars}
\end{figure*}

Gemini 지휘관을 쓴 두 조건이 기준선 대비 유의한 개선을 보인다(각각 $+0.060$, $+0.062$).
반면 Qwen2.5 7B는 두 기동 조건 모두에서 기준선보다 \emph{유의하게 나쁘다}($-0.192$,
$-0.203$). 이 음의 결과는 \ref{sec:res_threshold}절에서 다룬다.

\paragraph{포메이션별 분해}
\label{sec:res_ceiling}
통합 수치는 포메이션에 따라 다르게 읽어야 한다(Table~\ref{tab:ceiling}). 집중 포메이션은
휴리스틱만으로 포획률 0.997에 도달하여 개선 방향으로는 어떤 방법도 변별되지 않고, 통합값은
이 구간이 3분의 1을 차지하여 희석된다. 양동과 파상만 보면 제안 체계의 개선은 각각
$+0.090$, $+0.103$ 으로 통합값 $+0.062$ 보다 크다. 반면 악화 방향으로는 이 구간도 변별력을
가져, Qwen2.5 7B 배정은 $-0.157$ ($[-0.323, -0.057]$)로 거의 완벽한 상황마저 악화시킨다.
즉 \ref{sec:res_threshold}절의 능력 임계는 가장 쉬운 포메이션에서도 관측된다.

% 자동 생성 — tools/make_paper_tables.py. 손으로 고치지 말 것.
\begin{table}[htbp]
\centering
\caption{포메이션별 천장 효과. 집중 포메이션은 baseline 이 이미 0.997 이라 개선 여지가 0.003 뿐이므로 이 포메이션 단독으로 결론을 내면 안 된다.}
\label{tab:ceiling}
\footnotesize
\resizebox{\ifdim\width>\linewidth\linewidth\else\width\fi}{!}{%
\begin{tabular}{lcccc}
\toprule
포메이션 & baseline & 남은 여유 & 제안 & $\Delta$ \\
\midrule
집중 & 0.997 & 0.003 & 0.990 & -0.007 \\
양동 & 0.787 & 0.213 & 0.877 & +0.090$^{*}$ \\
파상 & 0.750 & 0.250 & 0.853 & +0.103 \\
\bottomrule
\end{tabular}
}
\\[3pt]
\begin{minipage}{\linewidth}\footnotesize
$^{*}$ 신뢰구간이 0 을 포함하지 않음.
\end{minipage}
\end{table}


\subsubsection{지휘관 능력 임계}
\label{sec:res_threshold}

기동 계층을 휴리스틱으로 고정하여 지휘관 축만 분리하면, 모델 능력에 따른 단조 증가가
나타난다(Fig.~\ref{fig:ladder}).

\begin{figure}[!tbp]
\centering
\includegraphics[width=0.85\columnwidth]{fig19.png}
\caption{지휘관 능력 사다리(기동 계층 휴리스틱 고정, 점선: 휴리스틱 배정 기준선)}
\label{fig:ladder}
\end{figure}

3개 포메이션 통합 기준으로 Qwen2.5 7B는 0.652, Qwen2.5 14B는 0.811, Gemini는 0.904이며,
휴리스틱 기준선은 0.844이다. 즉 \textbf{휴리스틱 선을 14B와 Gemini 사이에서 교차한다.}

함의는 ``LLM 지휘관을 쓰면 좋다''가 아니라 \textbf{``일정 능력에 못 미치는 지휘관은 코드
휴리스틱보다 못하다''}이며, 7B의 유의한 악화($-0.192$, $\dz = -1.00$)가 이 임계 주장의 근거다.

7B가 어디서 지는지는 휴리스틱 배정과의 \emph{일치율}이 보여 준다. 계획 호출마다 같은 상태에
\ref{sec:heuristic}절의 코드가 낼 배정을 함께 계산하고, 방어정별로 언어모델의 배정이 그것과 같은
비율을 재면(Fig.~\ref{fig:agreement}) 평균 일치율은 7B 0.391, 14B 0.538, Gemini 0.591 이고 세 방어정이
모두 일치한 호출의 비율은 각각 0.215, 0.365, 0.374 이다. 7B는 휴리스틱과 \emph{다르게} 배정하는
일이 가장 많고 그 결과는 기준선 아래다. 즉 임계 아래 모델의 이탈은 판단이 아니라 오류다. 다만
Gemini 도 호출의 41\,\%에서 휴리스틱과 다르게 배정하면서 기준선을 넘었으므로 일치율이 높을수록
좋다는 뜻은 아니며, 임계 위의 모델은 벗어나되 맞게 벗어난다.

\begin{figure*}[!tbp]
\centering
\includegraphics[width=0.7\textwidth]{fig20.png}
\caption{휴리스틱 배정과의 일치율. (a) 모델별 평균 일치율(95\,\% 신뢰구간)과 완전 일치 호출 비율(빗금),
         (b) 일치율 대 포획률(기동 계층 휴리스틱 고정; 파선은 휴리스틱 배정 기준선 0.844).}
\label{fig:agreement}
\end{figure*}

\subsubsection{지휘관 성능의 결정 요인}
\label{sec:res_mechanism}

능력 순위는 모델 세대가 바뀌면 의미를 잃으므로, 더 오래 남는 질문은 \emph{무엇이} 성능을 만드는가이다.
계획 단위 계측 결과를 Table~\ref{tab:planquality}와 Fig.~\ref{fig:planquality}에 정리하였다(점
하나가 지휘관 모델 하나, 오차막대는 양방향 부트스트랩 95\,\% 신뢰구간; $n=3$ 이라 회귀선·
상관계수는 얹지 않았다).

\begin{figure*}[!tbp]
\centering
\includegraphics[width=0.75\textwidth]{fig21.png}
\caption{계획 품질 대 성능: (a) 배정 커버리지, (b) 재계획 간 churn, (c) 경로 교차}
\label{fig:planquality}
\end{figure*}

% 자동 생성 — tools/make_paper_tables.py. 손으로 고치지 말 것.
\begin{table*}[t]
\centering
\caption{지휘관 계획 품질과 성능. 계획 품질은 시뮬레이션 결과와 독립적으로, 계획이 산출된 시점에 측정한다. \textbf{7B 는 커버리지가 가장 높은데 포획률은 가장 낮다}. 실패 원인은 미배정이 아니다.}
\label{tab:planquality}
\footnotesize
\setlength{\tabcolsep}{4pt}%
\resizebox{\ifdim\width>\linewidth\linewidth\else\width\fi}{!}{%
\begin{tabular}{lcccc}
\toprule
지휘관 & 커버리지 $\uparrow$ & churn $\downarrow$ & 교차 $\downarrow$ & 포획률 $\uparrow$ \\
\midrule
Qwen2.5 7B & 0.934 & 0.324 & 0.152 & 0.652 \\
Qwen2.5 14B & 0.877 & 0.167 & 0.111 & 0.811 \\
Gemini 3.5 Flash-Lite & 0.881 & 0.098 & 0.033 & 0.904 \\
\bottomrule
\end{tabular}
}
\\[3pt]
\begin{minipage}{\linewidth}\footnotesize
커버리지 $=$ 배정된 활성 클러스터 비율. churn $=$ 재계획 간 담당이 바뀐 배의 비율. 교차 $=$ 두 방어정의 요격 경로가 교차한 비율. 포획률은 기동 계층을 휴리스틱으로 고정한 조건(LLM 지휘관 + 휴리스틱 기동)의 값이다. 세 모델 모두 언어모델 호출이 실패해 휴리스틱이 대신 배정한 비율은 0.007 이하였으므로, 배정은 사실상 전부 언어모델이 낸 것이다.
\end{minipage}
\end{table*}


결정적 관측은 \textbf{Qwen2.5 7B가 커버리지는 가장 높은데(0.934) 포획률은 가장 낮다(0.652)}는
것이다. 즉 막아야 할 무리를 배정하지 않아서 지는 것이 아니다.

성능과 단조로 대응하는 것은 churn과 경로 교차이다. churn은 $0.324 \to 0.167 \to 0.098$,
경로 교차는 $0.152 \to 0.111 \to 0.033$ 으로 능력축을 따라 감소하며, 같은 순서로 포획률은
$0.652 \to 0.811 \to 0.904$ 로 증가한다(모두 기동 계층을 휴리스틱으로 고정한 조건). churn 0.324는 재계획마다 방어정의 약 3분의 1이
담당을 바꾼다는 뜻이며, 이는 \emph{반쯤 전개한 그물벽을 버리고 떠난다}는 것을 의미한다.
경로 교차는 두 방어정이 서로의 진로를 가로지르는 배정이다.

이 관찰이 제시하는 설계 원리는 \textbf{주기적으로 재계획하는 구조에서는 계획의 최적성보다
안정성이 중요하다}는 것이다(\ref{sec:discussion}절).

\subsection{두 기여의 결합}
\label{sec:res_g3}

\subsubsection{두 계층 기여의 가산성}
\label{sec:res_interaction}

식~\eqref{eq:interaction}의 상호작용을 10개 지표 전부에 대해 계산하였다
(Table~\ref{tab:interaction}, Fig.~\ref{fig:interaction}). 그림의 (a)에서 두 선이 평행에
가까울수록 두 계층의 기여가 가산적이고, (b)의 속이 빈 표식은 유의하지 않은
차이다. (a)에 오차막대를 두지 않은 것은 대응이 풀린 셀별 구간이 시나리오 난이도의 분산을
포함해 넓어 유의한 대응 차이를 가리기 때문이며, 불확실성은 전부 (b)에 있다.

% 자동 생성 — tools/make_paper_tables.py. 손으로 고치지 말 것.
\begin{table*}[t]
\centering
\caption{2$\times$2 상호작용 분해 (Gemini 지휘관). 10개 지표 중 \textbf{9개}에서 상호작용이 유의하지 않으므로 두 계층의 기여는 \textbf{가산적}이다.}
\label{tab:interaction}
\footnotesize
\resizebox{\ifdim\width>\linewidth\linewidth\else\width\fi}{!}{%
\begin{tabular}{lcc}
\toprule
지표 & 총효과 & 상호작용 [95\,\% CI] \\
\midrule
포획률 & +0.062 & -0.039 [-0.083, +0.012] \\
돌파 수 & -0.62 & +0.39 [-0.13, +0.83] \\
충돌률(척당) & -0.052 & -0.015 [-0.070, +0.022] \\
그물 걸림 & -0.60 & +0.17 [-0.08, +0.41] \\
포획당 그물 & -0.328 & +0.004 [-0.045, +0.059] \\
포획당 이동거리 (m) & -182 & -55 [-125, +30] \\
총 선회량 (rad) & -757 & -56 [-165, +59] \\
포획 이격거리 (m) & +223 & +38 [-2, +81] \\
평균 포획 시각 (step) & -22 & -12 [-20, -6]$^{*}$ \\
전개 시 적 거리 (m) & -371 & +26 [-44, +109] \\
\bottomrule
\end{tabular}
}
\\[3pt]
\begin{minipage}{\linewidth}\footnotesize
상호작용 $=$ (LLM+U-Net $-$ LLM+휴리스틱) $-$ (휴리스틱+U-Net $-$ baseline). 표식 $^{*}$ 은 유의한 상호작용이며, 나머지는 두 계층이 서로를 돕지도 잡아먹지도 않는다는 뜻이다.
\end{minipage}
\end{table*}


\begin{figure*}[!tbp]
\centering
\includegraphics[width=0.65\textwidth]{fig22.png}
\caption{$2\times2$ 상호작용(Gemini 지휘관): (a) 네 셀의 포획률 평균, (b) 시드 대응 차이와 BCa 신뢰구간}
\label{fig:interaction}
\end{figure*}

\textbf{10개 지표 중 9개에서 상호작용이 유의하지 않고, 다중비교 보정 후에도 그대로
9개다.} 두 계층은 서로를 돕지도 잠식하지도 않는다. 본 논문은
이를 상승효과로 해석하지 않고 \emph{가산적}이라고 보고한다. 두 계층이 서로 다른 실패
모드를 다루므로 기여가 독립적으로 누적된다.

한 가지 예외는 평균 포획 시각이다. 상호작용 $-11.98$, $95\,\%$ 구간 $[-19.84,\; -5.93]$ 의
유의한 상쇄이며 보정 구간 $[-23.31,\; -3.54]$ 에서도 살아남는다. 두 계층 모두 제압을 앞당기는
방향으로 작용하는데 \emph{같은 포획을 두 번 앞당길 수는 없으므로}, 공유 메커니즘에 대한
수확체감이다.

포메이션별로 보면 두 계층의 분업이 드러난다. 양동 포메이션에서는 기동 계층 단독
($-0.033$)도 전략 계층 단독($+0.070$)도 유의하지 않은데, 둘을
결합하면 $+0.090$ 으로 유의해진다. 무리 수와 방어정 수가 같아 여유가 없는 이 포메이션에서는
배정과 기동이 둘 다 맞아야 막히며, 2계층 설계의 효과가 드러나는 지점이 여기다. 반대로 단이
시간차를 두고 오는 파상에서는 기동 계층 단독($+0.157$)이 이미 유의하며 결합이 그보다 낫지
않다.

\subsubsection{지표별 효과크기와 다중비교 보정}
\label{sec:res_effectsize}

\ref{sec:metrics}절에서 정의한 10개 지표 전부에 대해 제안 체계와 기준선의 대응 차이를
계산하였다. 비교는 제안 체계 \texttt{llm\_unet}@Gemini 와 기준선 \texttt{heur\_heur} 사이이며,
결과는 Table~\ref{tab:effectsize}와 Fig.~\ref{fig:forest}에 있다. 그림의 $\dz$ 구간은 평균차의
BCa 구간을 $\mathrm{sd}(d)$ 로 나눈 것이라 표의 판정과 같은 사건이다.

% 자동 생성 — tools/make_paper_tables.py. 손으로 고치지 말 것.
\begin{table*}[tp]
\centering
\caption{전 지표 효과크기 (제안 = LLM 지휘관 + U-Net 점수맵, Gemini). $|d_z|$ 내림차순. \textbf{포획률이 가장 둔한 지표이며, 다중비교 보정에서 유의성을 잃는 유일한 방어 성과 지표다.}}
\label{tab:effectsize}
\footnotesize
\resizebox{\ifdim\width>\linewidth\linewidth\else\width\fi}{!}{%
\begin{tabular}{lccccc}
\toprule
지표 & 방향 & baseline & 제안 & $\Delta$ & $d_z$ \\
\midrule
총 선회량 (rad) & $\downarrow$ & 1625 & 868 & -757$^{*}$$^{\dagger}$ & -3.23 \\
포획당 그물 & $\downarrow$ & 0.737 & 0.409 & -0.328$^{*}$$^{\dagger}$ & -2.09 \\
포획 이격거리 (m) & $\uparrow$ & 1511 & 1734 & +223$^{*}$$^{\dagger}$ & +1.81 \\
전개 시 적 거리 (m) & $\downarrow$ & 1499 & 1128 & -371$^{*}$$^{\dagger}$ & -1.60 \\
그물 걸림 & $\downarrow$ & 0.82 & 0.22 & -0.60$^{*}$$^{\dagger}$ & -1.19 \\
평균 포획 시각 (step) & $\downarrow$ & 430 & 408 & -22$^{*}$$^{\dagger}$ & -1.04 \\
포획당 이동거리 (m) & $\downarrow$ & 825 & 643 & -182$^{*}$$^{\dagger}$ & -0.65 \\
충돌률(척당) & $\downarrow$ & 0.130 & 0.078 & -0.052$^{*}$$^{\dagger}$ & -0.52 \\
돌파 수 & $\downarrow$ & 1.56 & 0.93 & -0.62$^{*}$ & -0.44 \\
포획률 & $\uparrow$ & 0.844 & 0.907 & +0.062$^{*}$ & +0.44 \\
\bottomrule
\end{tabular}
}
\\[3pt]
\begin{minipage}{\linewidth}\footnotesize
방향 $\uparrow$ = 높을수록 좋음, $\downarrow$ = 낮을수록 좋음. $^{*}$ 95\,\% 신뢰구간이 0 을 포함하지 않음. $^{\dagger}$ 지표 10개에 대한 Bonferroni 보정(99.5\,\% 구간)에서도 0 을 배제함(단검이 없는 행은 보정을 견디지 못한다).
\end{minipage}
\end{table*}


\begin{figure}[!tbp]
\centering
\includegraphics[width=\columnwidth]{fig23.png}
\caption{제안 대 기준선의 전 지표 효과크기 $\dz$. 굵은 선은 신뢰구간, 옅은 선은 다중비교
         보정 구간이며, 속이 빈 표식은 보정 후 유의성을 잃는 지표다.}
\label{fig:forest}
\end{figure}

포획률의 효과크기는 $\dz = 0.44$ 로 측정한 10개 지표 중 최하위권이며, 총 선회량의 7분의 1
수준이다. 가장 크게 개선된 것은 기동 비용과 자원 효율이다. 총 선회량은 $1625$ 에서 $868$ 로
47\,\% 줄어 $\dz = -3.23$, 포획당 그물은 $0.737$ 에서 $0.409$ 로 45\,\% 줄어 $-2.09$,
그물 걸림은 $0.82$ 에서 $0.22$ 로 73\,\% 줄어 $-1.19$ 이고, 포획 이격거리는 $1511$ 에서
$1734\unit{m}$ 로 늘어 $+1.81$ 이다.

다중비교 보정을 걸면 이 대조가 더 선명해진다. 10개 중
8개는 그대로 유의하지만 포획률과 돌파 수는 유의성을 잃는다. 포획률의 보정 구간은
$[-0.003, +0.132]$ 로 유의성을 잃으며, \ref{sec:res_maneuver}절의 기동 계층 단독 조건도 같다.
보정을 견디지 못하는 것은 ``더 많이 막는다''는 주장이고, 견디는 것은 ``같은 정도로 막으면서
자원과 기동과 위험을 크게 줄인다''는 주장이다. 포획률의 효과크기가 작은 까닭은
\ref{sec:res_ceiling}절의 천장이며, 효율 지표에는 그러한 천장이 없다.

\subsection{지휘관 응답 지연}
\label{sec:res_latency}

응답 지연은 모델 규모에 단조가 아니다(Table~\ref{tab:latency_tail}). 로컬 14B가 16.59초인 데
비해 클라우드 모델은 1.56초로, \textbf{로컬이 10.6배 느리다}. 한 결정 주기가 25초
(25\,step $\times$ $\Delta t\!=\!1\unit{s}$)이므로 14B의 평균 지연은 주기의 \textbf{약
3분의 2}를 소모한다.

평균만 보면 주기 안에 들어오지만, 결정적인 것은 분포의 꼬리다
(Table~\ref{tab:latency_tail}, Fig.~\ref{fig:latency_ecdf}). 14B 호출 5{,}366건 가운데 \textbf{4.2\,\%가 25초 주기를
초과}하였고, 99 백분위는 27.8초, 최댓값은 296.4초였다. 동기 구조라면 이 4.2\,\%마다 제어
루프가 실제로 정지한다. 반면 클라우드 모델은 최댓값이 2.75초로 주기의 11\,\%에 그쳐
동기 호출로도 여유가 있다.

\begin{figure}[!tbp]
\centering
\includegraphics[width=\columnwidth]{fig24.png}
\caption{지휘관 응답 지연의 경험 누적분포(파선은 결정 주기)}
\label{fig:latency_ecdf}
\end{figure}

\begin{table*}[tp]
\centering
\caption{지휘관 응답 지연 분포. 결정 주기 25초와 비교한다. 로컬 14B만 주기를 초과하는
         꼬리를 갖는다.}
\label{tab:latency_tail}
\footnotesize
\begin{tabular}{lrrrrr}
\toprule
지휘관 & $n$ & 평균 & $p_{99}$ & 최대 & $>25$초 \\
\midrule
Qwen2.5 7B (로컬)    & 4{,}299 & 2.75  & 7.08  & 10.06  & 0.0\,\% \\
Qwen2.5 14B (로컬)   & 5{,}366 & 16.59 & 27.81 & 296.41 & \textbf{4.2\,\%} \\
Gemini (클라우드)    & 3{,}882 & 1.56  & 2.30  & 2.75   & 0.0\,\% \\
\bottomrule
\end{tabular}
\\[3pt]
\begin{minipage}{\linewidth}\footnotesize
단위는 초. 로컬 모델은 학습·평가와 GPU 한 대를 공유한 워크스테이션에서 4비트 양자화로 서빙하였다. 14B 의 최댓값 296초는 가중치 재적재를
시사하며, 이 꼬리가 로컬 중형 모델의 본질인지 자원 공유의 결과인지는 본 실험으로 구분되지
않는다. 읽어야 할 것은 로컬 중형 모델의 지연이 결정 주기를 넘길 수 있다는 정성적 사실이다.
\end{minipage}
\end{table*}

이 수치의 측정 구성에는 교란 요인이 있으나(Table~\ref{tab:latency_tail} 주석), 결론은
같다. \ref{sec:arch}절의 비동기 설계는 클라우드 배치에서는 선택이지만 로컬 배치에서는
지연 꼬리에 대한 안전장치다.

